% Auto-generated by extract_papers.py. Edit papers.json or the paper titleInput.tex files.
\section*{Research in Progress}
\textbf{Internal and External R\&D: An analysis of costs and benefits}\\[0.35em]
This paper analyzes the costs and benefits of internal and external R\&D activities. Using Dutch production and innovation surveys between 2000 and 2020, focusing on the ICT industry, I document an increasing share of firms performing R\&D and the four-way mix of in-house research, contracted-out research, both, and neither. To rationalize these findings, I build and estimate a dynamic discrete-choice model of R\&D with mode-specific investment costs. The cost of contracted-out R\&D is much higher than the cost of in-house R\&D, which explains the small share of firms that contract out alone. Doing both is cheaper than the sum of the two costs: in-house R\&D offsets almost all of the extra cost of going outside. I read that extra cost as a hold-up problem with the partner. The offset is absorptive capacity: a firm that already does in-house R\&D keeps more of the collaboration, so the two activities cost less together than apart. Productivity gains are similar across modes, so the mix of activities is driven by these costs. A uniform subsidy that does not favor a mode raises the share of R\&D-active firms and firm value by more than a subsidy aimed only at in-house R\&D.\par
\vspace*{0.05in}
\textit{Presentation}: \href{http://asesec.org/jornadas_economia_industrial/programme.html}{JEI - Gran Canaria} (2022), \href{https://sites.google.com/view/phd-evs2020/home?authuser=0}{Ph.D.-EVS} (2022), IO/Applied Micro Brownbag - Penn State (2023), VU Amsterdam Internal Seminar (2023), \href{https://earie.org/programme-rome-2023/}{EARIE - Rome} (2023), Tinbergen Institute Ph.D. Seminar (2023), DKIIS - Valencia (2023), \href{https://econ.la.psu.edu/comparative-analysis-of-enterprise-data-caed/}{CAED} - Penn State (2024), \href{https://economicdynamics.org/sedam_2024/}{SED Meeting} - Barcelona (2024), \href{https://us10.forward-to-friend.com/forward/preview?u=bb253d28cb05d3719c9e0c960&id=c939ef0aa0}{CEPR Workshop: Trade, Geography, and Industrial Organization - Erasmus Rotterdam} (2024, poster), \href{https://sites.google.com/view/apioc-2024}{APIOC - Seoul} (2024), CEPR Symposium - Paris (2024, poster), EEA-ESEM - Bordeaux (2025), CEPR IO Virtual Gathering (2025)\par
\vspace*{0.12in}
\textbf{Entry into Innovation Areas}\\[0.35em]
This paper investigates the multi-area entry decisions of publicly listed firms in the United States electronics industry from 1990 to 2019. I apply Latent Dirichlet Allocation to over 630,000 patent titles and abstracts to endogenously define twenty innovation areas and merge these with market-implied patent valuations. The data reveal two empirical patterns. First, innovation occurs across multiple domains simultaneously. Second, the number of active competitors varies across areas. To rationalize these patterns, I estimate a static simultaneous-move game of firm entry with moment inequalities. First, I find that firms enter areas where expected patent valuations are higher. Second, holding those valuations fixed, additional rivals reduce expected profit. Third, firms also benefit from occupying more areas at once. Fourth, a large per-area cost of occupying an area in a given year, together with rival congestion, still limits how many areas a firm enters. Starting from the areas firms already occupy, and holding rival counts fixed, I find that a 25\% reduction in the yearly per-area cost raises to 22--27\% the share of firms that occupy an extra area in a given year. A 50\% reduction in how costly rivals are raises that share to 15--21\%. In a case study of targeted cost reductions, those partial-equilibrium reallocations raise occupancy in semiconductor areas such as ``Semiconductor Memory Architecture'' and ``Semiconductor Lithography and Fabrication Processes.'' The offset is lower occupancy in other fields, including ``Wireless Base-Station and Mobile Communications.''\par
\vspace*{0.05in}
\textit{Presentation}: VU Amsterdam Internal Seminar (2025)\par
\vspace*{0.12in}
\textbf{The Impact of Government Financial Assistance on the Scope and Direction of Firm Innovation}, with C. Cincotta.\\[0.35em]
Governments spend hundreds of billions of dollars each year to support private-sector innovation, but it is unclear whether this money expands existing research, redirects it toward new technologies, or reallocates innovation away from other firms. To address this question, we develop a two-stage model in which firms first choose an R\&D portfolio and then compete in prices \`a la Salop. Financially constrained firms underinvest, and even unconstrained firms under-provide quality because product-market competition leaves a gap between private and social returns. Liquidity assistance can relax the cash constraint, while the quality gap remains. Cheaper cash can expand how much a recipient patents and how many technology areas it covers. If the firm simply scales up the research it already conducts, the composition of that research does not change. Assistance that is targeted toward particular activities can redirect the portfolio. Nearby rivals may expand through knowledge pooling or contract through business stealing. We take these implications to a panel of U.S.\ public firms from 2000 to 2021, combining Good Jobs First awards with patent data. Scope measures how much a firm patents and how many technology areas it covers. Direction measures whether those patents align with frontier or peer technologies. To estimate causal effects, we first reweight untreated firms so that they match recipients on pre-treatment characteristics, and we then estimate a stacked difference-in-differences. Our model implies that distant rivals provide a better control group for the treated firm's own response. Nearby rivals may themselves be affected by business stealing or knowledge pooling, so comparing the two groups traces how those nearby rivals respond. We implement this design using both industry codes and text-based product-market overlap. Distant rivals provide the preferred reading of the treated firm's own response. Financial assistance expands the scale of recipients' patent portfolios. Relative to distant industry controls, intensive scope rises by about 10.5 percent after an award. Relative to firms with little product-market overlap, the same contrast is about 3 percent. Direction shifts toward the industry's most-cited technologies, and the shift is strongest for cash grants. When nearby means high text-based product-market overlap and distant means little or none, the product-market spills are indistinguishable from zero, so those nearby rivals' patenting scale does not move. When nearby means the same two-digit industry, the close estimate on scale is smaller than the distant estimate, which means industry-code neighbors expand as well, the knowledge-pooling case.\par
\vspace*{0.12in}
\textbf{Mergers and Innovation: An Exploration of Scope and Direction of Innovation}, with S. Dobbelaere and J. L. Moraga-Gonzalez.\\[0.35em]
How do mergers change the scale and composition of firm innovation? We study U.S.\ public mergers from 1980 to 2020 using combined acquirer--target patent portfolios. We measure \emph{scope}, the scale of the joint portfolio, and \emph{direction}, its alignment with highly cited patents. A compact framework isolates three mechanisms. Common ownership internalizes business-stealing externalities in product-market overlap areas and reduces the private return to innovation there. The merger may also pool technological assets, raising the productivity of R\&D where the firms possess related knowledge. Some innovation lines require a fixed cost before they can be operated; pooled assets can make previously inactive lines worth activating. Product-market overlap and technological overlap are distinct: firms may compete for the same customers without sharing knowledge, or share related technologies without competing for the same customers. We compare treated merger pairs to matched control pairs before and after the deal. On the matched sample, joint-portfolio scope is larger after mergers than for matched controls. The difference is substantially larger when pre-merger technologies overlap, especially among non-horizontal deals. We find no strong evidence that activity is reallocated between technology classes shared by both firms and classes unique to one. Joint portfolios show higher alignment with highly cited patents relative to matched controls, although the framework leaves the sign of that change theoretically ambiguous. Acquirer-only R\&D measures overstate post-merger input growth relative to joint-entity measurement.\par
\vspace*{0.05in}
\textit{Presentation}: VU Amsterdam Internal Seminar (2024), Tinbergen Institute Workshop in Economics of Innovation (2025)\par
